\PassOptionsToPackage{dvipsnames,svgnames,table}{xcolor}
\documentclass[11pt]{article}
\usepackage[a4paper,margin=22mm]{geometry}
\usepackage[no-math]{fontspec}
\setmainfont{DejaVu Serif}
\usepackage{amsmath,amssymb,amsthm,mathtools,graphicx,xcolor,hyperref,xurl,xspace,ifthen,enumerate}
\usepackage[shortlabels]{enumitem}
\setlength{\emergencystretch}{4em}
\AtBeginDocument{\special{pdf:mapline pzdr Dingbats <uzdr.pfb}}
\SetMathAlphabet{\mathbf}{normal}{OT1}{cmr}{bx}{n}
\SetMathAlphabet{\mathbf}{bold}{OT1}{cmr}{bx}{n}
\newtheorem{theo}{Theorem}[section]
\newtheorem{lemm}[theo]{Lemma}
\newtheorem{thm}[theo]{Theorem}
\newtheorem{proposition}[theo]{Proposition}
\newtheorem{corollary}[theo]{Corollary}
\newtheorem{definition}[theo]{Definition}
\newtheorem{theorem}[theo]{Theorem}
\newtheorem*{remas*}{Remarks}
\newtheorem{remas}[theo]{Remarks}
\newtheorem*{exam*}{Example}
\newtheorem{ques}[theo]{Question}
\providecommand{\og}{«\nobreakspace}
\providecommand{\fg}{\nobreakspace»}
\newtheorem{lem}[theo]{Lemma}
\newtheorem{lemma}[theo]{Lemma}
\newtheorem{prop}[theo]{Proposition}
\newtheorem{coro}[theo]{Corollary}
\newtheorem{corol}[theo]{Corollary}
\newtheorem*{conj*}{Conjecture}
\newtheorem{conj}[theo]{Conjecture}
\newtheorem{conjecture}[theo]{Conjecture}
\newtheorem{Proposition}[theo]{Proposition}
\newtheorem{theoreme}[theo]{Theorem}
\theoremstyle{definition}
\newtheorem{defi}[theo]{Definition}
\newtheorem{exam}[theo]{Example}
\newtheorem{exem}[theo]{Example}
\newtheorem{remark}[theo]{Remark}
\newtheorem{example}[theo]{Example}
\newtheorem{exams}[theo]{Examples}
\newtheorem{rema}[theo]{Remark}
\newtheorem{nota}[theo]{Notation}
\newtheorem{notation}[theo]{Notation}
\newtheorem*{notas}{Notations}
\newtheorem{result}[theo]{Result}
\newtheorem{question}[theo]{Question}
\newtheorem{prob}[theo]{Problem}
\newtheorem*{rema*}{Remark}
\newtheorem*{defi*}{Definition}
\newtheorem*{theo*}{Theorem}
\newtheorem*{lemm*}{Lemma}
\newtheorem*{prop*}{Proposition}
\newcommand{\enoncename}{}
\newtheorem{corpusenonce}[theo]{\enoncename}
\newenvironment{enonce}[2][]{\renewcommand{\enoncename}{#2}\begin{corpusenonce}}{\end{corpusenonce}}
\newenvironment{enonce*}[2][]{\par\medskip\noindent\textbf{#2.}\itshape}{\par\medskip}
\newenvironment{proofc}[1][\proofname]{\begin{proof}[#1]}{\end{proof}}
\newcommand{\Mc}{\mathcal{M}}
\newcommand{\pointir}{.\ ---\ }
\newcommand{\equalenv}[2]{\expandafter\let\csname#1\expandafter\endcsname\csname#2\endcsname\expandafter\let\csname end#1\expandafter\endcsname\csname end#2\endcsname}
\providecommand{\diagbox}[3][]{\begin{tikzpicture}[x=1em,y=1em,baseline=(current bounding box.center)]\useasboundingbox(0,0) rectangle(3,1.4);\draw(0,1.4)--(3,0);\node at(.5,.3){#2};\node at(2.5,1.1){#3};\end{tikzpicture}}
\providecommand{\eg}{e.g.\ }
\providecommand{\ie}{i.e.\ }
\providecommand{\cf}{cf.\ }
\providecommand{\longthanks}[1]{\par\smallskip\noindent #1\par\smallskip}
\providecommand{\qbinom}[2]{\genfrac{[}{]}{0pt}{}{#1}{#2}}

\providecommand{\mathof}[1]{\mathrm{#1}}

\numberwithin{equation}{section}


\def\endappendix{}


\usepackage{latexsym,amsmath,amssymb,verbatim,tikz,tikz-cd}
\usepackage{mathtools}
\usepackage{graphicx,xcolor}

\usepackage[english]{babel}
\usepackage{amssymb}
\usepackage{amsmath}
\usepackage{amsthm}
\usepackage{graphicx}
\usepackage{csquotes}
\usepackage{bm}
\usepackage{enumitem}
\usepackage{subcaption}
\usepackage{caption}
\renewcommand{\thesubfigure}{\Alph{subfigure}}
\usepackage{dirtytalk}
\providecommand{\nopunct}{\spacefactor1007\relax}
\newtheorem{problem}[theo]{Problem}

\newcommand{\sym}{\mathrm{Sym}}
\newcommand{\qsym}{\mathrm{QSym}}
\newcommand{\Des}{\operatorname{Des}}
\newcommand{\Asc}{\operatorname{Asc}}
\newcommand{\Syt}{\operatorname{SYT}}
\newcommand{\Short}{\operatorname{Short}}
\newcommand{\Loop}{\operatorname{Loop}}
\newcommand{\Stable}{\operatorname{Stable}}
\newcommand{\len}{\operatorname{len}}
\newcommand{\height}{\operatorname{height}}
\newcommand{\core}{\operatorname{core}}
\newcommand{\res}{\operatorname{res}}
\newcommand{\row}{\operatorname{row}}
\newcommand{\insrt}{\operatorname{insert}}
\newcommand{\UP}{\operatorname{UP}}
\newcommand{\DOWN}{\operatorname{DOWN}}


\newcommand{\ZZ}{\mathbb{Z}}
\newcommand{\NN}{\mathbb{N}}
\newcommand{\RR}{\mathbb{R}}
\newcommand{\QQ}{\mathbb{Q}}
\newcommand{\PP}{\mathbb{P}}

\newcommand{\F}{\mathcal{F}}
\newcommand{\Q}{\mathcal{Q}}
\newcommand{\A}{\mathcal{A}}
\newcommand{\M}{\mathcal{M}}
\newcommand{\I}{\mathcal{I}}
\newcommand{\PM}{\mathcal{PM}}
\newcommand{\Pcal}{\mathcal{P}}


\newcommand{\odds}{\mathrm{odds}}

\makeatletter
\newcommand{\proofstep}[1]{
  \par
  \addvspace{\medskipamount}
  \textbf{#1\@addpunct{.}}\enspace\ignorespaces
}

\newcommand\restr[2]{{
  \left.\kern-\nulldelimiterspace 
  #1 
  \vphantom{\big|} 
  \right|_{#2} 
  }}
\makeatother












\title{Removing the first descent from a two-row standard Young tableau is a bijection}
\author{Avichai Marmor}
\date{doi:10.5802/alco.351; CC BY 4.0}
\begin{document}
\maketitle
\noindent\textit{Editorial note.} The counted statement is the author's labelled lemma `lem:remove first descent from syt` (source lines 459-466); the proof environment begins on the next line after \textbackslash{}end\{lemma\}, so the pairing is by adjacency.\par\medskip
\noindent\textit{Editorial note.} The retained context is the author's own lead-in sentence, which records that this observation is a crucial step in the inductive proof that follows.\par\medskip
\noindent\textit{Editorial note.} The retained preamble is the editorially prepared standalone preamble retained for this source, which already carries the author's full theorem-environment family and macro set; the publisher layout adapter of the source is not used, and the packages needed only by other parts of the article are deliberately left out because no retained line uses them.\par\medskip
\noindent\textit{Editorial note.} No retained line refers to a label, citation or figure outside the retained span, so the package declares no external reference.\par\medskip
\noindent\textit{Editorial scope.} One original result of the pinned publisher article, counted once: the shift-and-insert bijection for two-row standard Young tableaux (author label lem:remove first descent from syt).\par\medskip
\noindent\textit{Allowed standard prerequisites (editorial).}\begin{itemize}
\item Standard Young tableaux of shape (k\_1,k\_2). Syt(k\_1,k\_2) is the set of standard Young tableaux with two rows of lengths k\_1 >= k\_2, filled with 1,...,k\_1+k\_2 increasingly along rows and columns.. Both sides of the displayed identity are sets of such tableaux, and the map is defined on entries.
\item Descent set. Des(T) is the set of i such that i+1 appears in a strictly lower row than i in T.. The condition Des(T) contains J is what is transported by the bijection.
\end{itemize}\medskip
We start by introducing an observation, which will serve as a crucial step in our inductive argument.

\begin{lemma}
\label{lem:remove first descent from syt}
    Let $k_1 \ge k_2$ be two nonnegative integers, and let $J \subseteq [k_1 + k_2 - 1]$ be a set. Then
    \[
        \left|\Syt(k_1, k_2)(\Des \supseteq J)\right| = \left|\Syt\big(k_1+1, k_2+1\big)\big(\Des \supseteq (\{1\}\cup (J+2))\big)\right|,
    \]
    where $J+2 = \{j + 2 \mid j \in J\}$.
\end{lemma}

\begin{proof}
    The proof is bijective: We associate a given tableau $T \in \Syt(k_1, k_2)$ such that $\Des(T) \supseteq J$ with the tableau $T' \in \Syt(k_1+1, k_2+1)$ defined by $T'_{1,1} = 1,\; T'_{2,1} = 2$, and for all $i \in \{1,2\}$ and $j \in [k_i]$, $T'_{i,j+1} = T_{i,j}+2$ (recall that $T_{i,j}$ is the entry in row $i$ and column $j$ of $T$). Intuitively, we first increase every entry of $T$ by $2$. Then, we insert a new column with the letters $1,2$ on the left of $T$, shifting the other columns to the right. Clearly, this process is injective, and $T'$ satisfies $T' \in \Syt(k_1+1, k_2+1)$ and $\Des(T) \supseteq \{1\}\cup (J+2)$. Furthermore, any $T' \in \Syt(k_1+1, k_2+1)$ with $\Des(T') \supseteq \{1\}\cup (J+2)$ can be obtained through this process.
\end{proof}

\end{document}
